%---------------------------------------------------------------------------------
%                西南交通大学研究生学位论文：结论
%---------------------------------------------------------------------------------
\stepcounter{chapter}
\chapter*{结论与展望}
\addcontentsline{toc}{chapter}{结论与展望}
\setcounter{section}{0}

\section{主要研究结论}

本文围绕政务场景中自然语言问数与数据分析的实际需求，面向“如何可靠查数、如何规范分析、如何系统落地”三个核心问题开展研究，形成了从方法创新到系统落地的完整技术链路。全文主要研究结论如下。

在方法层面，本文提出两项核心创新。
第一，围绕政务术语歧义明显、统计口径隐含、复杂查询长尾分布突出等问题，本文提出了基于逻辑增强与异构协同的政务数据查询生成方法。该方法通过构建领域知识库、动态样例库与LA-Schema，将业务术语映射、逻辑计算公式和值域约束显式融入查询生成过程，并结合逻辑映射生成器、渐进分治生成器、执行反馈修复与候选判别模型，提升了复杂政务场景下查询生成的正确性与稳定性。实验结果表明，本文方法在自建政务数据集GovSQL上取得了91.47\%的执行准确率，较最强基线XiYan-SQL提升7.64个百分点；在BIRD与Spider公共数据集上分别达到71.23\%和89.25\%，表明所提出方法在垂直场景和通用场景中均具有较强竞争力。

第二，围绕开放式分析需求约束隐含、查询与分析职责耦合以及结果表达需规范化等问题，本文提出了基于分析规划与查询增强的政务数据分析方法。该方法引入AP-Schema作为结构化分析规划模式，将自然语言分析需求显式映射为分析类型、指标、维度、筛选条件、粒度、查询层操作、分析层操作、可视化规格和摘要要求等字段；进一步通过查询增强的任务分解机制，将涉及统计口径、字段绑定和复杂关联的操作下沉至查询层完成，在可靠取数基础上由轻量分析算子链与可视化规则生成规范化结果。实验结果表明，在政务分析任务上，本文方法的流程成功率、任务完成率、结果正确率和图表恰当性分别达到92.9\%、85.7\%、95.8\%和4.6分，显著优于代码生成（Code-Gen Agent）和工具迭代（ReAct Agent）等主流分析范式；在规划质量评估中，AP-Schema的意图识别准确率、槽位填充完备率和约束满足率分别达到94.6\%、91.8\%和91.1\%。这说明“先查准、再分析”的协同机制能够有效提升政务数据分析的稳定性与规范性。

在系统实现层面，本文完成了从方法到系统的工程整合。

第三，在前述查询方法和分析方法基础上，本文设计并实现了一个面向政务场景的智能问数系统。系统采用表示层、应用层、智能服务层和数据层四层架构，基于FastAPI、LangGraph、PostgreSQL、Milvus和Redis等技术组件，支持同步查询、异步分析、任务状态管理、执行日志留痕与结果追溯。系统实现表明，本文提出的查询方法与分析方法不仅能够在实验环境中取得有效结果，而且能够被组织为可交互、可复核、可追踪的统一系统，从而完成从自然语言需求理解、可靠取数到规范分析输出的闭环落地。

总体而言，本文构建了面向政务智能问数的完整研究框架，为大语言模型在复杂业务场景中的数据查询、分析与系统落地提供了可参考的方法路径。

\section{研究不足}

本文仍存在以下不足。

第一，数据与任务验证范围仍然有限。本文构建的GovSQL数据集及分析任务主要来源于同源政务业务场景，虽然能够较好反映财政预算、公共资源和政务项目等典型需求，但对于跨地区、跨部门、跨制度口径的复杂业务场景，其泛化能力仍有待进一步验证。

第二，领域知识库构建、动态样例整理和部分训练标注仍依赖人工参与。当前方法虽然通过结构化知识注入显著提升了垂直场景性能，但知识获取、校验与维护成本相对较高，面对业务规则频繁调整、数据结构持续变化的政务环境时，自动更新与自适应能力仍显不足。

第三，本文分析任务的覆盖范围仍以趋势、对比、排名和结构四类典型任务为主。对于更复杂的预测分析、诊断归因、多轮深度分析以及更严格的安全权限约束控制，本文尚未展开系统研究，相关能力距离真实复杂政务应用场景仍存在一定差距。

\section{未来展望}

第一，构建更大规模、多区域、多领域的政务查询与分析评测数据集，覆盖更多制度口径与数据库形态，进一步提升方法的跨场景泛化验证能力，并为后续模型比较与持续优化提供更充分的评测基础。

第二，研究更自动化的领域知识获取、知识更新与样例演化机制。可结合元数据同步、历史查询日志挖掘、人机协同校验与持续学习等方式，降低领域知识库与样例库的人工维护成本，增强系统对动态业务规则和模式变化的适应能力。

第三，拓展分析任务类型与系统能力边界。在方法层面，可进一步增强多轮交互分析、复杂洞察生成与多源异构数据协同分析能力；在系统层面，可进一步完善权限控制、部署治理与服务稳定性设计，推动大语言模型问数系统在真实政务场景中的安全应用与持续落地。
